\documentclass[11pt,a4paper]{article}

% ============================================================
%  Por qué "sin SSL" y el supervisado completo se mueven junto con los SSL
%  (split externo de 5 pliegues, protocolo v2; 2026-11-21).
%  Figuras, tablas y números (docs/figures_senal_compartida/):
%    python scripts/4_evaluate/trivial_baselines.py --check-init
%    python scripts/5_reports/report_shared_signal.py
%  Compilar desde docs/: pdflatex informe_senal_compartida.tex (dos veces)
% ============================================================

\usepackage[a4paper,margin=2.1cm]{geometry}
\usepackage[T1]{fontenc}
\usepackage{lmodern}
\usepackage[spanish,es-noshorthands]{babel}
\usepackage{microtype}
\usepackage{xcolor}
\usepackage{amsmath,amssymb}
\usepackage{booktabs}
\usepackage{array}
\usepackage{tabularx}
\usepackage{graphicx}
\usepackage{caption}
\usepackage{float}
\usepackage{placeins}
\usepackage{enumitem}
\usepackage{fancyhdr}
\usepackage{titlesec}
\usepackage[most]{tcolorbox}

\definecolor{primaryblue}{HTML}{1B4F72}
\definecolor{steelblue}{HTML}{2E86AB}
\definecolor{accentorange}{HTML}{E67E22}
\definecolor{softgray}{HTML}{85929E}
\definecolor{panelbg}{HTML}{EAF2F8}
\definecolor{warnbg}{HTML}{FDEBD0}
\definecolor{tldrbg}{HTML}{EBF5FB}
\definecolor{okc}{HTML}{1E8449}
\definecolor{noc}{HTML}{B03A2E}

\usepackage{xurl}
\usepackage[colorlinks=true,linkcolor=primaryblue,citecolor=primaryblue,urlcolor=steelblue]{hyperref}

\newtcolorbox{tldrbox}[1][En una frase]{colback=tldrbg,colframe=primaryblue,coltitle=white,
  colbacktitle=primaryblue,fonttitle=\bfseries,boxrule=0.9pt,arc=3pt,
  left=8pt,right=8pt,top=5pt,bottom=5pt,title=#1}
\newtcolorbox{ideabox}[1][Idea clave]{colback=panelbg,colframe=steelblue,coltitle=white,
  colbacktitle=steelblue,fonttitle=\bfseries,boxrule=0.8pt,arc=2pt,
  left=8pt,right=8pt,top=5pt,bottom=5pt,title=#1}
\newtcolorbox{warnbox}[1][Ojo con esto]{colback=warnbg,colframe=accentorange,coltitle=white,
  colbacktitle=accentorange,fonttitle=\bfseries,boxrule=0.8pt,arc=2pt,
  left=8pt,right=8pt,top=5pt,bottom=5pt,title=#1}

\pagestyle{fancy}
\fancyhf{}
\renewcommand{\headrulewidth}{0.5pt}
\fancyhead[L]{\small\color{softgray}Señal compartida · 5 pliegues externos · protocolo v2}
\fancyhead[R]{\small\color{softgray}SSL, sin SSL y supervisado}
\fancyfoot[C]{\small\thepage}
\setlength{\headheight}{14pt}

\titleformat{\section}{\normalfont\Large\bfseries\color{primaryblue}}{\thesection}{0.6em}{}
\titleformat{\subsection}{\normalfont\large\bfseries\color{steelblue}}{\thesubsection}{0.6em}{}

\graphicspath{{figures_senal_compartida/}}
\setlength{\parskip}{4pt}
\setlength{\parindent}{0pt}
\setlength{\emergencystretch}{3em}
\setlist[itemize]{topsep=2pt,itemsep=2pt}
\setlist[enumerate]{topsep=2pt,itemsep=2pt}
\captionsetup{font=small,labelfont={bf,color=primaryblue}}
\renewcommand{\arraystretch}{1.15}
\newcommand{\si}{\textcolor{okc}{\textbf{sí}}}
\newcommand{\no}{\textcolor{noc}{\textbf{no}}}

\input{figures_senal_compartida/numeros.tex}

\begin{document}

% ============================================================
\begin{center}
{\LARGE\bfseries\color{primaryblue} ¿Por qué \emph{sin SSL} y el supervisado\\[2pt]
se mueven junto con los SSL?}\\[8pt]
{\large\color{steelblue} La señal que comparten todos los modelos\\
cubos hipocampales de OASIS3 · 5 pliegues externos · protocolo v2}\\[6pt]
{\small 21 de noviembre de 2026}
\end{center}

\begin{tldrbox}[Respuesta corta]
Todos los modelos miden, sobre todo, \textbf{cuánto se atrofió el hipocampo}. Esa
señal es tan fuerte y tan simple que la recoge casi cualquier cosa que mire el cubo.
\begin{itemize}
  \item El \textbf{volumen hipocampal de FreeSurfer usado tal cual} (sin red, sin probe,
        sin entrenar nada) da AUROC de test \nvolTest. BYOL, el mejor SSL, da \nbyolTest.
  \item \textbf{38 estadísticas de intensidad del cubo}, con el mismo probe, dan \nstatsTest.
  \item Por eso los modelos suben y bajan juntos de un pliegue a otro: lo que cambia es
        qué tan separables son los sujetos que cayeron en cada test. Incluso el volumen,
        que no tiene ningún modelo detrás, sigue a la media de los SSL ($r = \nrVol$).
  \item \emph{Sin SSL} no es un caso raro: es una medición más ruidosa de la misma
        atrofia. Con pesos aleatorios la misma red baja a \nrandCV{} en CV, así que
        ImageNet sí aporta. El SSL suma unos $+0.03$ más y deja los embeddings más
        alineados con el volumen.
  \item Su ``mejora'' entre informes (de 0.698 a 0.803 en el pliegue 0) vino del
        arreglo del probe en el protocolo v2, no del encoder.
  \item No encontré fugas.
\end{itemize}
\end{tldrbox}

% ============================================================
\section{Qué se observa y qué se preguntó}

En la figura final (\texttt{docs/figures\_final/metricas.png}) quedan todos cerca: los
cuatro SSL (VICReg, BYOL, DINOv2 y rotación), el supervisado completo y el encoder
\emph{sin SSL}, que es el mismo ConvNeXt-T 3D con pesos ImageNet inflados, congelado,
más un probe lineal. El AUROC de test va de \ninitTest{} (sin SSL) a \nbyolTest{} (BYOL),
y los puntos de cada pliegue suben y bajan juntos. La pregunta es por qué, y sobre todo
por qué en el caso sin SSL, que nunca vio una resonancia durante su entrenamiento.

``Moverse al unísono'' puede leerse de tres maneras, y las tres tienen la misma causa de
fondo:
\begin{enumerate}
  \item \textbf{Entre pliegues}: un pliegue fácil lo es para todos (sección~\ref{sec:pliegues}).
  \item \textbf{De la CV al test}: en un mismo pliegue, el test queda por encima o por
        debajo de la CV en casi todos los modelos a la vez (sección~\ref{sec:pliegues}).
  \item \textbf{Entre informes}: sin SSL ``mejoró'' de 0.698 a 0.803 en el pliegue 0
        cuando cambió el protocolo (sección~\ref{sec:sinssl}).
\end{enumerate}

\section{Qué se comparó}

Todo se evaluó con el \textbf{mismo protocolo v2} (\texttt{t5\_ssl.evaluation.protocol\_v2})
y sobre los \textbf{mismos sujetos} en cada uno de los 5 pliegues externos:
\begin{itemize}
  \item unos 300 sujetos etiquetados y emparejados por edad para entrenar el probe;
  \item una CV representativa de unos 700 sujetos (23\,\% AD, edades reales);
  \item un test de unos 176 sujetos.
\end{itemize}
Los modelos son ensambles de 3 semillas, salvo sin SSL, que no depende de la semilla.

Además de los seis modelos de la figura final, agregué líneas base que no aprenden
ninguna representación (\texttt{scripts/4\_evaluate/trivial\_baselines.py}):
\begin{itemize}
  \item \textbf{Estadísticas de intensidad del cubo}. Por hemisferio: la media, la
        desviación, 7 cuantiles y un histograma de 10 tramos de las intensidades
        normalizadas del cubo $32^3$ que ven las redes. Son 38 números.
  \item \textbf{Volúmenes FreeSurfer}. Son los vóxeles de la parcelación
        \texttt{aparc+aseg} que ya trae el LMDB: hipocampo, amígdala, asta temporal,
        ventrículo lateral, corteza entorrinal y parahipocampal (cada una a ambos lados),
        más sustancia blanca, corteza y cerebro. Son 15 números. También probé solo los
        dos volúmenes hipocampales.
  \item \textbf{ConvNeXt-T con pesos aleatorios} (dos semillas), congelado: la misma
        arquitectura que sin SSL, pero sin ImageNet.
  \item \textbf{Un solo número como puntaje, sin probe}: el volumen hipocampal (menor
        volumen = AD) y la edad (mayor edad = AD).
\end{itemize}

Ninguna de estas features usa etiquetas. Las 38 estadísticas y los volúmenes pasan por
el mismo probe que los embeddings de 768 dimensiones: regresión logística balanceada,
con C elegido por log-loss en la CV interna. De los \nnPooled{} sujetos de test, \nnVol{}
tienen parcelación FreeSurfer. A los que no la tienen, el probe les imputa la media de
los etiquetados del pliegue; en los análisis por sujeto se los excluye.

% ============================================================
\section{Resultado 1: una señal simple alcanza el techo}
\label{sec:escalera}

\begin{table}[H]
\centering\footnotesize
\setlength{\tabcolsep}{4pt}
\input{figures_senal_compartida/tab_resumen.tex}
\caption{Media de los 5 pliegues externos.
\emph{CV}: CV representativa (23\,\% AD, edades reales).
\emph{CV emp.}: la misma CV restringida a los sujetos etiquetados, que están emparejados
por edad; ahí la edad sola da \nageMatched.
\emph{Compuesto}: $(\text{AUROC}+\text{AUPR}+\text{F1-macro})/3$.
\emph{AUROC $\pm$3 años}: AUROC sobre los pares AD-HC cuya edad difiere en 3 años o menos.
Los puntajes sin probe no tienen umbral, así que tampoco tienen F1 ni compuesto.}
\label{tab:resumen}
\end{table}

\begin{figure}[H]
\centering
\includegraphics[width=\textwidth]{escalera.png}
\caption{Cada fila es un modelo o una línea base. El punto es la media de los 5 pliegues
y la barra va del peor al mejor pliegue. Los rombos son puntajes sin probe; los círculos
blancos, líneas base que pasan por el probe.}
\label{fig:escalera}
\end{figure}

La tabla~\ref{tab:resumen} y la figura~\ref{fig:escalera} ordenan las señales de menor a
mayor. En la CV representativa queda así:
\begin{center}\small
\begin{tabular}{@{}l c@{}}
pesos aleatorios & \nrandCV \\
sin SSL (ImageNet) & \ninitCV \\
SSL & 0.836--0.851 \\
estadísticas de intensidad & \nstatsCV \\
volumen hipocampal solo & \nvolCV \\
15 volúmenes FreeSurfer & \nasegCV \\
\end{tabular}
\end{center}
En el test, el volumen solo (\nvolTest), las 38 estadísticas (\nstatsTest) y BYOL
(\nbyolTest) quedan empatados. Tres cosas saltan a la vista:

\begin{enumerate}
  \item \textbf{El techo es la volumetría.} Un número que FreeSurfer calcula sin ver
        ninguna etiqueta alcanza al mejor SSL. Donde la edad no ayuda, lo supera: en la
        CV emparejada, el volumen hipocampal da \nvolMatched{} y los 15 volúmenes
        \nasegMatched, frente a \nmodelsMatchedMin--\nmodelsMatchedMax{} de los seis
        modelos. Los 15 volúmenes también tienen el mejor compuesto de test
        (\nasegComposite, frente a como mucho \nmodelsCompositeMax{} de los modelos).
  \item \textbf{La señal está en la intensidad misma del cubo.} Cuando el hipocampo se
        atrofia, el cubo tiene más líquido (oscuro) y menos tejido. 38 números que
        resumen el histograma de intensidades bastan para llegar al nivel de los SSL.
        Una red con pesos aleatorios, que solo combina intensidades locales al azar, ya
        saca \nrandCV. (El probe con los dos volúmenes hipocampales rinde algo menos que
        el volumen solo, en parte porque a los sujetos sin parcelación les imputa la media.)
  \item \textbf{Todos los modelos viven en una franja estrecha}: la de una medición
        más o menos limpia de la atrofia. El piso razonable no es 0.5 sino 0.75 (pesos
        aleatorios), y el techo está en 0.86--0.87 (volumetría).
\end{enumerate}

\subsection*{¿Y la edad?}
En la población real, la edad sola da \nageTest{} de AUROC, porque los AD son mayores.
Parte de lo que miden los modelos es atrofia por edad. Pero no es lo que los hace
moverse juntos:
\begin{itemize}
  \item Entre pares AD-HC con edades a 3 años o menos, los modelos siguen en
        \ninitAge--\nbyolAge{} y el volumen en \nvolAge, mientras la edad cae a 0.55.
  \item Los pliegues en que la edad separa mejor no son los mismos en que los modelos
        rinden más: la correlación con la media SSL es $r = \nrAge$
        (tabla~\ref{tab:pliegues}).
\end{itemize}

% ============================================================
\section{Resultado 2: por qué suben y bajan juntos entre pliegues}
\label{sec:pliegues}

\begin{figure}[H]
\centering
\includegraphics[width=\textwidth]{pliegues.png}
\caption{AUROC por pliegue. La banda azul es el rango de los cuatro SSL y la línea azul,
su media. La línea negra discontinua es el volumen hipocampal usado tal cual: no tiene
modelo, ni probe, ni semillas.}
\label{fig:pliegues}
\end{figure}

\begin{table}[H]
\centering\footnotesize
\setlength{\tabcolsep}{4pt}
\input{figures_senal_compartida/tab_pliegues.tex}
\caption{AUROC de test por pliegue y correlación de Pearson de cada fila con la media de
los cuatro SSL a través de los 5 pliegues. Abajo, la brecha test $-$ CV de cada pliegue.}
\label{tab:pliegues}
\end{table}

La figura~\ref{fig:pliegues} es la clave. La línea negra discontinua no es un modelo:
es el volumen hipocampal de cada sujeto de test, usado tal cual. Aun así sube y baja
exactamente con los SSL ($r = \nrVol$ a través de los pliegues) y con el supervisado
($r = \nrSup$ entre supervisado y media SSL). El pliegue 3 es fácil para todos y el 1
difícil para todos porque \textbf{en el pliegue 3 tocaron sujetos cuya atrofia separa
mejor a los AD de los HC}: el volumen solo da 0.893 ahí. En el pliegue 1, en cambio, los
volúmenes de AD y HC se solapan más (0.825). No es algo que hagan los modelos.

Cuantitativamente, entre los seis modelos ConvNeXt-T el pliegue explica el
\textbf{\nshareFold} de la varianza del AUROC de test; el modelo, el \nshareModel; y la
interacción, el \nshareInter. El test de cada pliegue tiene unos 176 sujetos (41 AD) y
su AUROC varía $\pm 0.03$ solo por qué sujetos tocaron. Las diferencias entre SSL
($\le 0.015$ en la media) son más chicas que eso.

\textbf{La brecha CV $\to$ test tiene la misma causa.} La CV de cada pliegue tiene unos
700 sujetos y casi no se mueve; el test, con 176, sí. La brecha test $-$ CV de la media
SSL (+0.010, $-$0.031, +0.009, +0.059, +0.029) sigue a la del volumen sin modelo
($-$0.001, $-$0.048, $-$0.018, +0.039, +0.028). Cuando un test ``sale alto'' para todos,
sale alto también para una regla que no se entrenó en nada.

\textbf{Sin SSL es el que menos acompaña} ($r = \nrInit$). Cae a 0.80 en los pliegues 0
y 4, donde los demás no caen. Es lo esperable de una medición más ruidosa de la misma
cantidad (sección~\ref{sec:sinssl}).

% ============================================================
\section{Resultado 3: sujeto a sujeto, ordenan casi igual}
\label{sec:sujetos}

\begin{table}[H]
\centering\footnotesize
\setlength{\tabcolsep}{4pt}
\input{figures_senal_compartida/tab_sujetos.tex}
\caption{Test agrupado de los 5 pliegues (\nnVol{} sujetos con volumen FreeSurfer). Cada
puntaje se pasa a percentil dentro de su pliegue.
\emph{AUROC con volumen similar}: AUROC sobre los pares AD-HC cuyo volumen hipocampal
difiere en 0.25 desviaciones estándar o menos; mide lo que el puntaje sabe más allá del
volumen.
$R^2$: fracción de la varianza del percentil del puntaje que explica linealmente el
percentil del volumen; entre paréntesis, la del volumen más la edad.}
\label{tab:sujetos}
\end{table}

\begin{figure}[H]
\centering
\includegraphics[width=\textwidth]{sujetos.png}
\caption{Cada punto es un sujeto del test agrupado, con puntajes en percentil dentro de
su pliegue. Izquierda: el puntaje de BYOL frente al volumen hipocampal (\nnVol{} sujetos
con volumen). Derecha: frente al puntaje sin SSL (los \nnPooled{} sujetos).}
\label{fig:sujetos}
\end{figure}

Si los modelos solo coincidieran en la etiqueta, sus puntajes se parecerían entre clases
pero no dentro de cada clase. No es lo que pasa:
\begin{itemize}
  \item \textbf{Entre modelos.} BYOL ordena a los sujetos casi igual que VICReg (0.95),
        el supervisado (\nrhoByolSup) y las 38 estadísticas de intensidad (0.93).
        También casi igual que sin SSL (\nrhoByolInit).
  \item \textbf{Dentro de cada clase.} Entre los AD solos, BYOL y sin SSL tienen
        Spearman \nrhoByolInitAD; entre los HC solos, \nrhoByolInitHC. Todos ordenan
        por la misma cantidad continua, el grado de atrofia, y no solo por la etiqueta.
  \item \textbf{Frente al volumen.} Con el volumen hipocampal, BYOL tiene Spearman
        \nrhoByolVolAD{} dentro de los AD y \nrhoByolVolHC{} dentro de los HC.
        \textbf{El volumen solo explica el \nrTwoByol{} de la varianza del puntaje BYOL}
        (\nrTwoByolAge{} con la edad); en sin SSL, el \nrTwoInit.
  \item \textbf{Más allá del volumen.} Entre pares AD-HC de volumen parecido, todos los
        modelos separan algo más: AUROC 0.62--0.66, \ncondByol{} en BYOL. Ese resto es
        igual de chico para todos, y las 38 estadísticas de intensidad también lo tienen
        (\ncondStats). La edad aporta poco ahí (\ncondAge).
\end{itemize}

% ============================================================
\section{El caso sin SSL}
\label{sec:sinssl}

\textbf{Es de verdad el encoder ImageNet congelado.} Recalculé sus embeddings desde cero:
la arquitectura de \texttt{train.py}, pesos ImageNet de torchvision inflados a 3D, sin
ningún paso de entrenamiento. Coinciden con los guardados por las corridas sin SSL en
los 5 pliegues, con un error relativo máximo de $1.4\times10^{-6}$ (redondeo de
float32). Al pasarlos por el protocolo dan exactamente los números guardados (pliegue 0:
CV 0.828, test 0.803).

\textbf{Por qué ya rinde 0.82.} El cubo trae la señal a nivel de intensidades. Los
filtros de ImageNet (bordes, manchas, texturas) inflados a 3D responden a los bordes
entre líquido y tejido. Sus 768 features contienen la atrofia, mezclada con muchas
direcciones irrelevantes. El probe lineal, con unos 300 sujetos y una regularización
fuerte (C entre 0.001 y 0.003 en casi todos los pliegues), la recupera solo en parte:
\ninitCV{} en CV, contra \nrandCV{} con pesos aleatorios.

\textbf{Qué agrega el SSL.} El SSL no descubre una señal nueva. Reorganiza las features
para que las direcciones en que más varían los cubos hipocampales queden explícitas, y
lo que más varía entre esos cubos es justamente la atrofia (edad más enfermedad). Por
eso:
\begin{itemize}
  \item la correlación con el volumen sube de \nrhoInitVol{} (sin SSL) a \nrhoByolVol{} (BYOL);
  \item el AUROC sube $+0.02$ a $+0.03$ en CV y $+0.02$ a $+0.04$ en test;
  \item el resto más allá del volumen apenas cambia (\ncondInit{} frente a \ncondByol).
\end{itemize}
El SSL le facilita el trabajo al probe, pero no lo lleva más allá del techo volumétrico.

\textbf{La ``mejora'' entre informes fue del probe, no del encoder.} En el pliegue 0, sin
SSL pasó de 0.698 a 0.803 de AUROC de test entre el protocolo v1 y el v2
(figura~\ref{fig:historia}). En el mismo cambio los SSL subieron solo $+0.015$ a
$+0.024$, y el supervisado quedó igual, porque sus predicciones salen de la red y no del
probe. El protocolo v1 elegía la C del probe de test con 60 sujetos de validación, en
una grilla de 0.01 a 10 que casi siempre tocaba su borde inferior. Así sobreajustaba, y
el más perjudicado era el encoder de 768 features más ruidosas: ImageNet. La CV
emparejada del v1 nunca mostró una ventaja grande del SSL (sin SSL 0.717, SSL
0.70--0.72; tabla~\ref{tab:historia}). La vieja ventaja de ``+0.12 a +0.16'' del SSL era
ese artefacto; la real es de unos $+0.03$.

\begin{figure}[H]
\centering
\includegraphics[width=0.7\textwidth]{historia.png}
\caption{Pliegue 0 (\texttt{split\_deploy}), media de las semillas de cada modelo,
evaluado con el probe del protocolo v1 y con el del v2. El supervisado no usa probe:
sus predicciones son las mismas en los dos protocolos.}
\label{fig:historia}
\end{figure}

\begin{table}[H]
\centering\small
\input{figures_senal_compartida/tab_historia.tex}
\caption{Pliegue 0 con los dos protocolos (media de semillas de corridas individuales,
no ensambles). En el v1, la CV era la emparejada por edad.}
\label{tab:historia}
\end{table}

% ============================================================
\section{Por qué el supervisado completo no supera al SSL}

El supervisado parte de los mismos pesos ImageNet y ve el mismo cubo. Se entrena de
punta a punta con unos 240 sujetos etiquetados y elige su época con 60 de validación.
Con tan pocas etiquetas, la red aprende lo más discriminativo y lo más fácil de aprender,
que es la atrofia. Llega al mismo techo (test \nsupTest, CV \nsupCV) y ordena a los
sujetos casi igual que BYOL (Spearman \nrhoByolSup).

Su CV es más baja y más variable (0.796 a 0.859 según el pliegue): cada pliegue de la CV
entrena una red nueva con unos 190 sujetos. Ningún método pasa el techo volumétrico
porque lo que queda más allá del volumen es poco (AUROC 0.62--0.66 entre pares de
volumen parecido) y repartido. Con este número de etiquetas, ni el SSL ni el supervisado
lo aprovechan mejor que 38 estadísticas de intensidad.

% ============================================================
\section{Auditoría de fugas}

Que un encoder congelado de ImageNet rinda 0.82 podría sugerir que algo se filtra. Revisé
cada punto por donde podría hacerlo:

\begin{table}[H]
\centering\small
\begin{tabularx}{\textwidth}{@{}X c@{}}
\toprule
Comprobación & ¿Sin fuga? \\
\midrule
Sujetos de test disjuntos de los de CV y pool en los 5 pliegues (0 en común; los 880
sujetos de test son todos distintos) & \si \\
Ninguna sesión de un sujeto de test entra al entrenamiento SSL (los manifiestos no
comparten sesiones; cada sujeto está en un solo split) & \si \\
Sin SSL es el encoder ImageNet congelado (embeddings recalculados = guardados, error
$\le 1.4\times10^{-6}$) & \si \\
Todos los modelos se evalúan en los mismos sujetos de test por pliegue & \si \\
Posición del cubo: centroide hipocampal medio de 429 sesiones de entrenamiento del
split~361, calculado solo con la anatomía (sin diagnóstico), y la misma para todas las
sesiones & \si \\
Normalización de intensidad por volumen (percentiles 1--99 de esa misma resonancia),
sin estadísticas de otros sujetos & \si \\
C del probe por log-loss en la CV interna; umbral de F1 elegido en la CV representativa,
nunca en el test & \si \\
Supervisado: entrena con train, elige época con val y usa el test solo para puntuar & \si \\
\bottomrule
\end{tabularx}
\end{table}

Sobre la posición del cubo: algunas de esas 429 sesiones caen en el test de algún pliegue
externo. Pero lo único que se toma de ellas es el promedio de la posición del hipocampo
en el espacio MNI, sin etiqueta, y es una sola posición fija para todos. No puede llevar
información de diagnóstico.

El argumento más fuerte, en todo caso, es otro: una regla que no se entrena en nada (el
volumen hipocampal) alcanza el mismo AUROC. No hace falta ninguna fuga para explicar
los números.

% ============================================================
\section{Qué significa para el proyecto}

\begin{ideabox}[Lo que cambia en la lectura de los resultados]
\begin{enumerate}
  \item \textbf{La línea base que hay que superar es la volumetría, no sin SSL.} Hoy
        ningún modelo supera al volumen hipocampal (test \nvolTest; CV emparejada
        \nvolMatched). Los 15 volúmenes FreeSurfer superan a todos los modelos en la CV
        emparejada (\nasegMatched) y en el compuesto de test (\nasegComposite). Conviene
        poner esta fila en cada tabla de resultados.
  \item \textbf{Las diferencias entre SSL están por debajo del ruido del test.} El test
        de cada pliegue se mueve $\pm 0.03$ por composición; las diferencias entre los
        cuatro SSL son de 0.015 o menos. Hay que compararlos con bootstrap pareado
        sobre los mismos sujetos o en el test agrupado de 880 sujetos, como ya se hace,
        y no por pliegue.
  \item \textbf{La tarea AD vs HC con el cubo hipocampal está saturada por la atrofia.}
        Para saber si el SSL aprende algo más, hacen falta preguntas donde el volumen
        no baste:
        \begin{itemize}
          \item medir lo que se sabe más allá del volumen (el AUROC entre pares de
                volumen similar de la tabla~\ref{tab:sujetos});
          \item curvas con menos etiquetas (10\,\%, 25\,\%), donde el SSL suele mostrar
                su ventaja;
          \item tareas más finas: CDR~0.5 frente a HC, o progresión;
          \item cohortes externas.
        \end{itemize}
  \item \textbf{Combinar en vez de competir.} Volumen y embedding se pueden juntar en un
        mismo probe. Para que cuente, la mejora tiene que validarse en la CV, no
        mirando el test.
\end{enumerate}
\end{ideabox}

\begin{warnbox}[Límites de este análisis]
\begin{itemize}
  \item Los volúmenes FreeSurfer se cuentan en la grilla MNI de 2\,mm: el registro al
        atlas ya corrige en buena medida el tamaño de cabeza, y la resolución es gruesa.
  \item 48 de los 880 sujetos de test no tienen parcelación; quedan fuera de los
        análisis por sujeto, y en el probe se les imputa la media.
  \item Las líneas base son una sola corrida por pliegue y los modelos, ensambles de 3
        semillas. Eso favorece levemente a los modelos.
  \item Las estadísticas de intensidad podrían recoger diferencias de escáner además de
        atrofia. Pero el volumen FreeSurfer, que no depende de la intensidad absoluta,
        da lo mismo, y los puntajes se ordenan como el volumen.
\end{itemize}
\end{warnbox}

\appendix
\section{Cómo se reproduce}

Todo corre en CPU, en unos 15 minutos, y no modifica ninguna corrida existente.
\begin{enumerate}
  \item \texttt{python scripts/4\_evaluate/trivial\_baselines.py -{}-check-init}\\
        Calcula las features de las 2059 sesiones del cache \texttt{data/roi3d/hippo\_s32}:
        estadísticas, ImageNet, pesos aleatorios y volúmenes del LMDB. Comprueba que los
        embeddings ImageNet sean los de las corridas sin SSL y evalúa cada conjunto con
        el protocolo v2 en los 5 pliegues. Escribe en
        \texttt{campaign\_ssl3d\_runpod/outer5/baselines/} los archivos
        \texttt{features.npz} y \texttt{protocol\_v2.json}.
  \item \texttt{python scripts/5\_reports/report\_shared\_signal.py}\\
        Lee además los ensambles de \texttt{report\_outer5.py} y escribe las figuras, las
        tablas y \texttt{numeros.tex} en \texttt{docs/figures\_senal\_compartida/}, y
        \texttt{summary.json} junto a las líneas base.
  \item Desde \texttt{docs/}: \texttt{pdflatex informe\_senal\_compartida.tex} (dos veces).
\end{enumerate}

\end{document}
